\documentclass[conference]{IEEEtran}

\usepackage[T1]{fontenc}
\usepackage[utf8]{inputenc}
\usepackage{cite}
\usepackage{amsmath,amssymb}
\usepackage{array}
\usepackage{booktabs}
\usepackage{graphicx}
\usepackage{multirow}
\usepackage[table]{xcolor}
\usepackage{tabularx}
\usepackage{hyperref}

\definecolor{TrustStackBlue}{HTML}{1F4E79}
\definecolor{TrustStackTeal}{HTML}{0F766E}
\definecolor{TrustStackGold}{HTML}{A56A00}
\definecolor{TrustStackSlate}{HTML}{334155}
\definecolor{TrustStackSoftBlue}{HTML}{F4F8FC}
\definecolor{TrustStackSoftTeal}{HTML}{F2FBFA}
\definecolor{TrustStackSoftGold}{HTML}{FFF8ED}

\hypersetup{
  colorlinks=true,
  linkcolor=TrustStackBlue,
  citecolor=TrustStackTeal,
  urlcolor=TrustStackGold
}

\newcommand{\TrustSection}[1]{\section{\texorpdfstring{\textcolor{TrustStackBlue}{#1}}{#1}}}
\newcommand{\TrustSubsection}[1]{\subsection{\texorpdfstring{\textcolor{TrustStackTeal}{#1}}{#1}}}
\newcommand{\TrustCallout}[2]{%
  \begin{center}
  \setlength{\fboxsep}{8pt}%
  \fcolorbox{TrustStackBlue!70!black}{TrustStackSoftBlue}{%
    \begin{minipage}{0.95\linewidth}
    \textbf{\textcolor{TrustStackBlue}{#1}}\par
    \smallskip
    #2
    \end{minipage}%
  }%
  \end{center}
}

\title{TrustStack: A Local-First, Evidence-Grounded Evaluation Standard for Trustworthy AI}

\author{
\IEEEauthorblockN{Anonymous Submission}
\IEEEauthorblockA{Prepared in IEEE conference format for TrustStack system review}
}

\begin{document}

\maketitle

\begin{abstract}
TrustStack is a local-first system for evaluating whether AI outputs are grounded, auditable, and safe enough to trust in operational settings. The system combines document ingestion, retrieval-augmented response generation, structured confidence scoring, risk labeling, operator-facing explanation, and report-artifact export into a single evaluation workflow. This paper introduces the \emph{TrustStack Evaluation Standard v2.0}, a weighted, evidence-first framework that scores answers across retrieval relevance, evidence sufficiency, citation traceability, claim support, contradiction risk, completeness, abstention behavior, answer discipline, safety risk, and calibration. We also describe an operator-facing review interface that exposes the full standardized suite, generates IEEE-ready export artifacts, and returns an aggregate score breakdown with reproducibility metadata. The result is a system intended not merely to answer questions, but to teach users how trust should be inspected before AI output is operationalized.
\end{abstract}

\begin{IEEEkeywords}
trustworthy AI, retrieval-augmented generation, explainability, evaluation standard, human-in-the-loop oversight, local-first systems
\end{IEEEkeywords}

\TrustSection{Introduction}
AI systems are increasingly deployed into settings where an answer is not enough; users must also understand why the answer should or should not be trusted. Existing chat interfaces typically optimize for fluency rather than auditability, leaving operators with limited visibility into evidence quality, citation traceability, contradiction risk, and the conditions under which a model should abstain. TrustStack addresses this gap by treating trust as an evaluation product rather than a by-product of generation.

TrustStack is motivated by one central question: \emph{can a user understand not only what the model said, but whether the answer is sufficiently grounded to act on?} To answer this, TrustStack combines local document ingestion, retrieval, answer generation, evaluation, and explanation into a single review surface. The system is explicitly designed for low-friction local deployment with MongoDB-backed persistence, live browser-based interaction, and deterministic fallback modes when external providers are unavailable.
The technical design builds on established retrieval-augmented generation and hallucination literature, while extending structured evaluation ideas into a local-first operator workflow \cite{lewis2020rag,gao2023ragsurvey,ji2023survey,openai2024evals}.

\TrustCallout{Paper Contributions}{
\begin{enumerate}
  \item It introduces the TrustStack Evaluation Standard v2.0, a weighted scoring framework for grounded-answer review.
  \item It describes a local-first implementation that unifies ingestion, retrieval, query, scoring, risk flags, and run history.
  \item It presents an operator-facing review workflow and export artifact path that map system behavior to report-ready outputs.
\end{enumerate}
}

\TrustSection{Problem Setting}
TrustStack assumes a corpus of uploaded evidence, a user question, and a model-generated answer conditioned on retrieved context. The system does not attempt to solve open-domain truth verification. Instead, it asks whether the answer is justified \emph{with respect to the indexed evidence currently available}. This is a narrower but more defensible setting for enterprise, academic, and operational review workflows.

The practical failure modes in this setting are familiar:
\begin{itemize}
  \item retrieval returns weak or irrelevant chunks;
  \item the answer cites evidence poorly or not at all;
  \item specific claims are unsupported by the retrieved material;
  \item the answer drifts beyond the scope of the question;
  \item the answer contains operational advice that requires human review;
  \item the interface fails to teach the user how the score was formed.
\end{itemize}

TrustStack is built to expose each of these failure modes rather than hide them behind a single scalar score.

\TrustSection{System Overview}
The TrustStack pipeline contains five operational layers:
\begin{enumerate}
  \item \textbf{Evidence Intake}: user uploads are parsed, chunked, and persisted.
  \item \textbf{Retrieval and Query}: the system embeds the question, retrieves the top evidence chunks, and assembles a grounded context.
  \item \textbf{Answer Generation}: a configured LLM or fallback path produces an answer, citations, and an insufficient-evidence signal.
  \item \textbf{Evaluation and Risk Analysis}: the answer is scored under TrustStack Evaluation Standard v2.0 and labeled with risk flags.
  \item \textbf{Operator Review}: the interface surfaces the answer, supporting chunks, diagnostics, historical runs, and standardized-suite summary.
\end{enumerate}

\begin{table}[t]
  \caption{TrustStack implementation stack}
  \label{tab:stack}
  \centering
  \rowcolors{2}{TrustStackSoftBlue}{white}
  \begin{tabular}{ll}
    \toprule
    \rowcolor{TrustStackBlue!12}
    \textcolor{TrustStackBlue}{Layer} & \textcolor{TrustStackBlue}{Implementation} \\
    \midrule
    Frontend & React, Vite, TypeScript \\
    Backend API & FastAPI \\
    Document store & MongoDB \\
    Retrieval store & ChromaDB or local vector fallback \\
    Embeddings & Ollama, local lexical fallback \\
    Generation & Ollama or extractive fallback \\
    Verification & unittest, Playwright, smoke tests \\
    \bottomrule
  \end{tabular}
\end{table}

\TrustSection{TrustStack Evaluation Standard v2.0}
The TrustStack Evaluation Standard v2.0 is a weighted framework that converts a single answer into an evidence-first evaluation report. Table~\ref{tab:dimensions} summarizes the ten dimensions currently implemented.
The report structure is aligned with structured evaluation practices in which each score is paired with a rationale, diagnostics, and a review-oriented interpretation \cite{openai2024evals}.

\begin{table*}[t]
  \caption{TrustStack Evaluation Standard v2.0}
  \label{tab:dimensions}
  \centering
  \rowcolors{2}{TrustStackSoftBlue}{white}
  \begin{tabular}{p{0.17\linewidth}p{0.08\linewidth}p{0.62\linewidth}}
    \toprule
    \rowcolor{TrustStackBlue!12}
    \textcolor{TrustStackBlue}{Dimension} & \textcolor{TrustStackBlue}{Weight} & \textcolor{TrustStackBlue}{Purpose} \\
    \midrule
    Retrieval relevance & 0.16 & Measures whether the retrieved evidence is strongly aligned with the question rather than dominated by one weak hit. \\
    Evidence sufficiency & 0.13 & Measures whether the answer is backed by enough non-duplicative support across the corpus. \\
    Citation traceability & 0.12 & Measures whether citations point back to the retrieved evidence and cover the answer transparently. \\
    Claim support & 0.15 & Measures whether individual claims in the answer are explicitly supported by retrieved evidence. \\
    Contradiction risk & 0.10 & Measures whether the answer appears to conflict with the retrieved evidence. \\
    Completeness & 0.09 & Measures whether the answer addresses the major terms and demands of the question. \\
    Honesty and abstention & 0.10 & Rewards the system for acknowledging missing evidence instead of overstating certainty. \\
    Answer discipline & 0.06 & Measures whether the answer stays proportional to the evidence without drifting into unsupported detail. \\
    Safety and operational risk & 0.05 & Measures whether the answer contains procedural or high-stakes guidance that needs review. \\
    Calibration and consistency & 0.04 & Measures whether the score and explanation are calibrated to detected weaknesses. \\
    \bottomrule
  \end{tabular}
\end{table*}

Each answer is also accompanied by:
\begin{itemize}
  \item explicit checks such as evidence presence, citation match, contradiction scan, question coverage, and safety screening;
  \item claim-level diagnostics identifying supported and weakly supported statements;
  \item evidence diagnostics including top-hit score, average hit score, supporting chunk count, source diversity, and unsupported claim ratio;
  \item recommended follow-up actions for the user.
\end{itemize}

\TrustCallout{Evaluation Output}{
The evaluation report is intentionally structured so that the same output can support live UI review, conference slides, and manuscript-ready artifact export. In practice, the report now behaves like a reusable evidence packet rather than a transient chat summary.
}

The score bands are currently defined as:
\begin{itemize}
  \item \textbf{Pass}: score $\geq 80$
  \item \textbf{Review}: score $\geq 60$ and $< 80$
  \item \textbf{Fail}: score $< 60$
\end{itemize}

\TrustSection{Standardized Test Suite}
In addition to single-query evaluation, TrustStack now runs a standardized suite that evaluates the current corpus and runtime stack through a set of grounded probes. The suite includes:
\begin{itemize}
  \item direct evidence retrieval prompts derived from the corpus;
  \item citation-audit prompts that ask which evidence supports a requirement;
  \item negative-control prompts that probe abstention on unsupported or out-of-scope questions;
  \item operational-restraint prompts that test whether high-stakes guidance is flagged for review;
  \item synthesis prompts that assess coverage across the uploaded material.
\end{itemize}

The suite aggregates per-answer evaluation reports into weighted category scores:
\begin{itemize}
  \item grounding and retrieval;
  \item auditability and traceability;
  \item safety and operational restraint;
  \item consistency and calibration;
  \item communication and explanation;
  \item corpus coverage.
\end{itemize}

\TrustCallout{Export and Report Artifact Path}{
TrustStack now produces a structured report artifact alongside the live score. The artifact contains the overall score, weighted category breakdown, case-level results, strengths, weaknesses, recommended follow-ups, and the metadata needed to regenerate IEEE-style tables and summary figures without manual transcription.
}

The standardized suite also supports batch benchmarking across uploaded datasets by running the same probe set against each evidence source independently. Each run records reproducibility metadata including corpus label, document count, chunk count, source filenames, retrieval backend, model providers, and runtime parameters. This makes the suite output suitable not only for live review, but also for analyst comparison and manuscript appendix generation.

\TrustSection{Operator Interface}
TrustStack exposes its evaluation workflow through an operator-facing review interface rather than a conventional chat-only surface. The interface is structured around the same stages used by the backend pipeline: evidence intake, runtime querying, evidence review, score interpretation, historical comparison, and export. This matters because the user is not merely consuming an answer. Instead, the user is asked to inspect whether a generated answer is sufficiently supported to justify operational use.

The interface is therefore designed around three principles. First, every trust judgment should remain traceable to retrieved evidence. Second, every score should be paired with diagnostics that explain what raised or lowered confidence. Third, evaluation outputs should be reusable beyond the immediate interaction, including benchmarking, analyst review, and manuscript export. In practice, this means the frontend behaves as an inspection layer over the standardized backend evaluation process rather than as a thin wrapper around generation alone.

For presentation and teaching use, the same interface can also be organized into a nine-part academic narrative: motivation, system architecture, runtime evaluation loop, evidence review, score interpretation, benchmarking, full-stack integration, methodology, and exportable artifacts. This ordering mirrors the structure of a conference-style systems talk and keeps the live demonstration aligned with the written report. The important point is not the visual arrangement itself, but the fact that the interface already exposes the underlying stages in a sequence that supports both operator review and formal explanation.

\TrustSection{Implementation and Verification}
TrustStack is built as a local-first system. MongoDB stores documents, chunks, and run history. The backend exposes ingestion, document listing, sample-question generation, query, run history, standardized-suite execution, batch benchmark, and report-export routes. The frontend consumes those routes through a single API layer and renders an operator-facing inspection workflow for evidence, evaluation, and reporting.

Verification currently operates at four levels:
\begin{enumerate}
  \item backend API and service unit tests;
  \item real MongoDB integration tests;
  \item browser E2E tests against the live app stack;
  \item served-stack smoke tests against built frontend assets, live FastAPI, and Docker MongoDB.
\end{enumerate}

This layered verification strategy matters because TrustStack is itself an evaluation product. If the system is intended to judge groundedness and traceability, its own ingest, retrieval, explanation, and persistence paths must be tested with equal rigor.

\TrustSection{Experimental Plan}
The current implementation enables several near-term experiments suitable for a conference evaluation section:
\begin{itemize}
  \item compare TrustStack score stability across repeated runs on the same corpus;
  \item compare grounded answers against an answer-only baseline with no retrieval or citation discipline;
  \item measure how often unsupported claims are detected when the question is partially or fully out of scope;
  \item measure operator-facing usefulness of diagnostics such as contradiction notes, unsupported claim ratio, and recommended actions;
  \item evaluate category-level sensitivity under intentionally degraded retrieval conditions.
\end{itemize}

The present codebase already supports deterministic local testing through lexical retrieval and fallback generation, which makes the system suitable for reproducible experiments without depending on cloud-only infrastructure.

\TrustSection{Limitations}
TrustStack currently evaluates groundedness relative to the uploaded corpus, not objective world truth. If the corpus is incomplete, outdated, or internally inconsistent, the system can still return a well-structured but ultimately limited evaluation. Likewise, the contradiction and claim-support checks are heuristic rather than formal entailment models. This is a pragmatic choice for local-first operation, but it limits semantic precision.

The current interface is still optimized for explanation and demonstration rather than dense analyst workflows, even though the backend now exports report-ready artifacts, per-case citation alignment metrics, batch benchmark summaries, and reproducibility metadata. Future work should therefore focus less on raw score export and more on deeper semantic entailment checks, denser analyst review tooling, and larger empirical evaluations across external benchmark corpora.

\TrustSection{Conclusion}
TrustStack reframes trustworthy AI interaction as an evidence-first evaluation problem. Instead of returning an answer alone, the system returns a score, a rationale, a set of diagnostics, and a review path. By combining a local-first architecture, a standardized evaluation framework, and an operator-facing inspection interface, TrustStack provides a concrete path toward AI systems that can be inspected before they are trusted.

\TrustSection{Acknowledgment}
This document is structured as a conference-style system paper scaffold. Final camera-ready submission details, author metadata, and empirical results tables should be updated once the full experimental runs are complete.

\bibliographystyle{IEEEtran}
\bibliography{references}

\end{document}
